\section{循环神经网络基础}

\subsection{RNN 的核心思想}

RNN 的核心特征是引入了一个\textbf{隐藏状态}（hidden state），该状态在处理序列时不断更新。每接收一个新输入，RNN 都会基于当前输入和上一步的隐藏状态，计算新的隐藏状态。

\begin{figure}[H]
\centering
\includegraphics[width=0.95\textwidth]{slides-images/slide-006.jpg}
\caption{RNN 的两种表示：左侧为折叠图（带自环箭头），右侧为展开图（explicitly showing 时间步之间的依赖关系）\protect\footnotemark}
\end{figure}
\footnotetext{Slides 第7页。}

\begin{importantbox}{RNN 的数学公式}
RNN 在每个时间步 $t$ 执行以下计算：

\textbf{隐藏状态更新}：
$$h_t = f_W(h_{t-1}, x_t)$$

\textbf{输出计算}：
$$y_t = f_Y(h_t)$$

其中：
\begin{itemize}
    \item $h_t$：时间步 $t$ 的隐藏状态向量
    \item $x_t$：时间步 $t$ 的输入向量
    \item $f_W$：带参数 $W$ 的状态转移函数（每个时间步共享相同的参数）
    \item $f_Y$：将隐藏状态映射到输出的函数
\end{itemize}
\end{importantbox}

\subsection{Vanilla RNN 的具体形式}

最常见的 RNN 形式（称为 Vanilla RNN）使用 $\tanh$ 作为激活函数：

$$h_t = \tanh(W_{hh} h_{t-1} + W_{xh} x_t + b)$$
$$y_t = W_{hy} h_t$$

\begin{figure}[H]
\centering
\includegraphics[width=0.95\textwidth]{slides-images/slide-009.jpg}
\caption{Vanilla RNN 的数学公式：隐藏状态通过 $\tanh$ 激活函数更新，输出通过线性变换从隐藏状态计算得到\protect\footnotemark}
\end{figure}
\footnotetext{Slides 第10页。}

\begin{knowledgebox}{为什么使用 $\tanh$ 激活函数}
$\tanh$ 的输出范围为 $[-1, 1]$，具有以下优势：
\begin{itemize}
    \item \textbf{有界性}：反复应用不会导致隐藏状态爆炸
    \item \textbf{零中心}：能表示正负值，比 sigmoid 更适合作为隐藏状态的激活函数
    \item \textbf{非线性}：提供必要的表达能力
\end{itemize}
但 $\tanh$ 的导数最大值为 1（在 $x=0$ 处），这意味着梯度在反向传播时几乎总是在衰减——这是导致梯度消失的重要原因之一。
\end{knowledgebox}

\subsection{三组权重矩阵}

RNN 中有三组关键的权重矩阵：

\begin{table}[H]
\centering
\begin{tabular}{lll}
\toprule
\textbf{权重矩阵} & \textbf{作用} & \textbf{维度} \\
\midrule
$W_{xh}$ & 将输入 $x_t$ 变换到隐藏状态维度 & $d_h \times d_x$ \\
$W_{hh}$ & 将上一步隐藏状态变换为当前步的贡献 & $d_h \times d_h$ \\
$W_{hy}$ & 将隐藏状态映射到输出 & $d_y \times d_h$ \\
\bottomrule
\end{tabular}
\caption{Vanilla RNN 的三组权重矩阵}
\end{table}

\begin{warningbox}{所有时间步共享同一组权重}
RNN 在每个时间步使用\textbf{完全相同}的权重矩阵 $W_{hh}$、$W_{xh}$、$W_{hy}$。这意味着：
\begin{enumerate}
    \item 模型大小不随序列长度增长
    \item 梯度计算时需要将各时间步的梯度累加
    \item 模型对序列中不同位置施加相同的变换规则
\end{enumerate}
\end{warningbox}

\subsection{本章小结}

RNN 通过引入隐藏状态和参数共享的循环结构，实现了对变长序列的建模。其核心计算包含两个矩阵乘法（输入到隐藏、隐藏到隐藏）加一个非线性激活函数，以及一个从隐藏状态到输出的线性映射。

%% ========== 手工构造 RNN ==========

